\subsection{Genus-zero faithfulness for the structured all-types action}

This statement concerns the labelled and coloured version of the
Hausdorff presentation in E.8. Compatibility means continuity, preservation
of types and of covering and filling generators, and compatibility with
the labelled/coloured finite-cover groupoids, their label-forgetting maps,
and all specified puncture-deletion maps. In particular it includes the
finite permutation quotient recording the ordering of any one colour of
punctures. This is the structure needed to state forgetting a chosen
subset on an unordered moduli stack. The theorem does not concern an
unspecified family of abstract automorphisms of unlabelled groups with no
such compatibility data.

\begin{claim}{E.13}{Proved}
Let $F$ be a continuous, type-preserving autoequivalence of the structured
completed category just specified. Suppose that its outer action on every
genus-zero mapping-class group, including its labelled and coloured
variants, is the identity. Then its outer action on every
$\widehat M_{g,n}$ is the identity.

Consequently restriction to the complete genus-zero outer-action system
is injective on the compatible systems of outer mapping-class-group
automorphisms induced by these autoequivalences. This asserts faithfulness
of the restriction on those systems. It does not by itself characterize
the arithmetic image, prove that every genus-zero automorphism lifts,
or identify autoequivalences that induce the same outer system on all
objects.
\end{claim}
\begin{proof}
We use markings to work with one object of each required marked
topological type. Changes of markings conjugate the group maps below.
All occurrences of an equality of fundamental-group homomorphisms are
understood up to one target inner automorphism, unless based paths have
been chosen. No assertion that the completed mapping-class group acts
faithfully on a profinite surface group is used.

\emph{An auxiliary continuous fundamental-group functor.}
For a surface $S$, put $P_S=\widehat{\pi_1(S)}$. An oriented covering
or filling induces an outer continuous homomorphism on these groups.
Diagram isotopy is a homotopy of the composite maps, and equal actual
composites induce equal outer homomorphisms. Hence these assignments
define a functor on $\mathcal C_{\rm diag}$.

For finitely generated profinite groups $P,Q$, the space
$\operatorname{Hom}_{\rm cont}(P,Q)$ is compact: using finitely many
generators of $P$, it is the closed subspace of a finite power of $Q$
cut out by all its profinite relations. Quotienting by the compact
conjugation action of $Q$ gives a compact Hausdorff space. Composition
is continuous, as can be checked in every finite characteristic
quotient of the target. The automorphism group
$\operatorname{Out}(P_S)$ is profinite. The universal property in E.6--E.8
therefore extends this functor to the completed category.

A completed covering arrow still gives an injection of index its degree:
it is an actual covering arrow composed with an element of
$\widehat M_S$, and the latter acts by a continuous automorphism of
$P_S$. A completed filling gives a surjection. These facts require the
existence of the mapping-class action on $P_S$, not its faithfulness.

Here is a useful consequence for completed diagrams. Suppose
\[
\begin{array}{ccc}
 S_D&\xrightarrow{\ f'\ }&T_q\\
 \downarrow u&&\downarrow v\\
 S&\xrightarrow{\ f\ }&T
\end{array}                                                     \tag{E.13.1}
\]
commutes, where $f,f'$ are degree-$d$ covers, $v$ fills one puncture
$q$, and $u$ fills $d$ punctures. Choose the target conjugator supplied
by the outer equality, and let $J'\subset P_{T_q}$ and $J\subset P_T$
be the images of the two covering maps. Since $u_*$ is surjective,
\[
 v_*(J')=J,
 \qquad J'\subset v_*^{-1}(J),
 \qquad [P_{T_q}:J']=[P_{T_q}:v_*^{-1}(J)]=d.
\]
Thus $J'=v_*^{-1}(J)$. It follows from the finite-cover classification
that $f'$ has precisely the right cover orbit of the puncture pullback
of $f$. In particular the inertia of $q$ acts trivially on its sheets.
Each of the $d$ source cusps deleted by $u$ maps to $q$: its image is
killed by $v_*$, whereas the inertia of every old target cusp survives
that filling. There are exactly $d$ degree-one cusps above $q$, so these
are all of them. This argument rules out a spurious completed square;
it does not assume a categorical pullback property for the isotopy
quotient.

\emph{Restricted cancellation after filling ordinary holes.}
Let $f:S\to T$ be a finite cover of affine hyperbolic surfaces and
$u:S_D\to S$ fill a finite set $D$ of additional ordinary punctures.
Retain $D$ as a colour, with further distinctions inside it if desired,
so the forgetting homomorphism $q:M_{S_D}\to M_S$ is defined on the
whole indicated group. Then the discrete right stabilizer is
\[
                    S_{fu}=q^{-1}(\operatorname{Deck}(f)).     \tag{E.13.1a}
\]
To prove the nontrivial inclusion, suppose $(fu)a=fu$ and put
$a_0=q(a)$. Apply the ordinary fundamental-group functor before
completion. Surjectivity of $u_*$ gives
\[
 f_*\,(a_0)_*=\operatorname{int}(t)\,f_*
                       \quad\text{for some }t\in\pi_1(T).
\]
Thus $t$ normalizes the finite-index subgroup
$J=f_*\pi_1(S)$. The corresponding element of
$N_{\pi_1(T)}(J)/J=\operatorname{Deck}(f)$ has the same outer action
on $\pi_1(S)$ as $a_0$. Since $S$ is aspherical, the two
homeomorphisms are homotopic, hence isotopic. The exact topological
input is \citet[Theorem 5.6(1); Corollary 5.7]{E_CantwellConlon2014}:
a finite-type hyperbolic surface without boundary has a complete
finite-area metric, admits no embedded hyperbolic half-plane, and has
only annular, not simply connected, ends. It satisfies those hypotheses.
This is the ordinary homotopy-to-isotopy theorem, equivalently the
injectivity part of the discrete Dehn--Nielsen--Baer theorem; no
profinite analogue is assumed. This proves (E.13.1a); the reverse
inclusion follows from $f a_0=f$ for a deck transformation.

The finite deck group embeds in the completed mapping-class group.
One direct check in this affine situation is as follows. It acts freely
on $S$, and $S/\operatorname{Deck}(f)$ is affine hyperbolic. Lifting
a finite graph spine of that quotient gives, over $\Q$, the module
\[
 H_1(S,\Q)\simeq\Q\oplus
 \Q[\operatorname{Deck}(f)]^{\,r-1},\qquad
 r=\operatorname{rank}\pi_1(S/\operatorname{Deck}(f))\ge2.
\]
This follows from the cellular chain sequence and semisimplicity, just
as in F.11. In particular the action is faithful and is detected modulo
some integer on integral homology. Thus the finite deck group survives
in a finite quotient of the mapping-class group.

The forgetting homomorphism $q$ is surjective. Right exactness of
completion gives $\ker\widehat q=\overline{\ker q}$. Since
$q^{-1}(\operatorname{Deck}(f))$ is a finite union of cosets of
$\ker q$, E.6 now gives the exact completed stabilizer
\[
 \operatorname{Stab}_{\widehat M_{S_D}}(fu)
       =\widehat q^{-1}(\operatorname{Deck}(f)).                \tag{E.13.1b}
\]
The formula holds in arbitrary completed markings as well. Indeed a
completed covering is $f_0\beta$, a completed filling is $u_0\gamma$,
and surjectivity of $\widehat q$ lifts $\beta$ through the source of
$u_0$. Their composite is therefore $f_0u_0\delta$ for a completed
source marking $\delta$. Transporting (E.13.1b) by this marking gives
the asserted formula with the forgetting map and deck group of the
given completed arrows. This supplies restricted cancellation on
covering arrows for every completed square, including its image under
$F$; it is not an epimorphism assertion for arbitrary filling arrows.

\emph{Recovering the surface group by one moving puncture.}
We only need this construction for affine hyperbolic surfaces, for which
the surface group is free of finite rank at least two. Let $S^+$ denote
$S$ with one additional distinguished puncture $p$, and let
$w:S^+\to S$ fill it. The discrete Birman exact sequence is
\[
 1\longrightarrow\pi_1(S,p)\longrightarrow M_{S^+}
 \xrightarrow{\operatorname{forget}_p}M_S\longrightarrow1.
\]
It is the fundamental-group sequence of the universal punctured curve
over the marked moduli stack. The family is locally topologically
trivial, by trivializing its compactification together with its disjoint
marked sections. The base stack is aspherical because its Teichmuller
space is contractible, as in E.8. The homotopy sequence therefore gives
the displayed injection and surjection.
The action of $M_{S^+}$ on the based group $\pi_1(S,p)$ extends to
an action on $P_S$. The completed free group $P_S$ is centre-free,
as proved in F.11. The composite
\[
 P_S\longrightarrow\widehat M_{S^+}
 \longrightarrow\operatorname{Aut}(P_S)
\]
is the inner-automorphism map, and hence is injective. Right exactness
of profinite completion then gives
\[
 1\longrightarrow P_S\longrightarrow\widehat M_{S^+}
 \longrightarrow\widehat M_S\longrightarrow1.                \tag{E.13.2}
\]
Equivalently $P_S$ is the right stabilizer of the filling arrow $w$.
This is a full profinite Birman sequence; it is not an identification
with a procongruence quotient.

The one-puncture filling has one right orbit. Its right stabilizer is
normal, so any compatible $F$ preserves this kernel and determines an
outer automorphism $\psi_S$ of $P_S$. Using the finite labelled
variants specifies which new puncture is being forgotten. If $c$ is an
old labelled puncture of $S$ and $S_c$ is still affine hyperbolic, the
commuting two-deletion square induces the usual surjection
\[
 P_S\longrightarrow P_{S_c}.
\]
Indeed a point-pushing loop is sent to the same loop after filling $c$;
the assertion follows first discretely and then by (E.13.2). Its kernel
is the closed normal subgroup
\[
 N_c=\overline{\langle\!\langle I_c\rangle\!\rangle}_{P_S}.
\]
Compatibility with the specified deletion maps gives
$\psi_S(N_c)=N_c$. No claim that a normal-closure kernel by itself
identifies a particular procyclic generator is needed.

\emph{Chosen-sheet naturality for a covering.}
For a connected finite cover $f:S\to T$, choose an ordinary point
$q\in T$, write $D=f^{-1}(q)$, and distinguish one point $p\in D$.
Use the square (E.13.1), with the distinguished point $p$ and the
remaining $d-1$ points assigned separate colours. Also use the filling
\[
 r:S_D\longrightarrow S_p
\]
of $D\setminus\{p\}$, followed by the one-point filling
$S_p\to S$. A cover automorphism fixing a point of the fibre is the
identity. Thus the right stabilizer of $f'$ in the mapping-class group
preserving the distinguished point $p$ is trivial, also after completion.
Its target lifting subgroup is open, and lifting gives a continuous
homomorphism on that subgroup, by E.6.

Let $P_T\subset\widehat M_{T_q}$ be the point-pushing subgroup in
(E.13.2). Inside it, consider the elements $b$ for which there is a lift
$a$ satisfying
\[
 b f'=f'a,
 \qquad u a=u.                                               \tag{E.13.3}
\]
The second condition says that, after the moving fibre is filled, the
source marking is the identity. It excludes adjusting a lift by a
nonidentity deck transformation. The first condition, together with the
distinguished point $p$, makes the lift unique. Discretely, the subgroup
of such $b$ is exactly
$f_*\pi_1(S,p)\subset\pi_1(T,q)$: a loop in $T$ returns the chosen
sheet precisely when its lift beginning at $p$ is closed. Forgetting the
other sheets by $r$ sends its lifted braid to that closed point-pushing
loop in $S_p$.

The same assertion holds for the full profinite groups. To check the
closure issue explicitly, intersect the open target lifting subgroup
with $P_T$. Since $v b=v$ for a target point-push, the lift relation
and the square give $(fu)a=fu$. Formula (E.13.1b) shows that the
forgotten source marking $a_0$ lies in the finite deck group of $f$.
This is valid for arbitrary completed squares, rather than only actual
cover transports. The discrete intersection with $P_T$ is dense in
the completed intersection, since the target lifting subgroup is open.
The condition $a_0=1$ in (E.13.3) is therefore the kernel of a finite
continuous quotient of that open subgroup. Its discrete subgroup is
dense, so it is the full completion of the just described index-$d$
subgroup. The map induced by $r$ identifies it continuously with $P_S$.
Under that identification its inclusion into $P_T$ is precisely $f_*$.

All ingredients in (E.13.3) are categorical equations, specified
fillings, and the finite cover that distinguishes one sheet. Their
images under $F$ have the same interpretation. The preceding
index calculation ensures that the image of (E.13.1) is the correct
puncture-pullback cover orbit, including its $d$ new degree-one cusps.
We obtain the naturality identity
\[
 (Ff)_*\,\psi_S=\psi_T\,f_*                                \tag{E.13.4}
\]
up to a target inner automorphism. The identity for completed arrows
follows by continuity on each right orbit. Thus the induced surface-group
automorphisms are natural for covers as well as for fillings. This
constructs the needed naturality; it is not an assumption that $F$
commutes with the auxiliary fundamental-group functor.

\emph{The genus-zero normalization.}
Choose $T$ of genus zero with $b\ge5$ punctures, with any prescribed
partition into labelled points and colours.
The groups $\widehat M_T$ and $\widehat M_{T_q}$ are centre-free.
For completeness, the pure genus-zero group with four marks is a free
group of rank two. Induction with (E.13.2) proves centrelessness of all
its higher pure completions. We also need the faithful action of the
permutation group $S_b$ on its rational first homology for $b\ge5$.
Write $e_{ij}=e_{ji}$ for the positive meridian of a collision of two
labelled punctures. Then
\[
 H_1(M_{0,b}^{\rm pure},\Q)
 =\Q^{\binom b2}/
       \langle\,\sum_{j\ne i}e_{ij}:1\le i\le b\,\rangle.   \tag{E.13.4a}
\]
Indeed, fixing three punctures at $0,1,\infty$ identifies the moduli
space with the complement of the hyperplanes $z_i=0$, $z_i=1$ and
$z_i=z_j$ in $\C^{b-3}$. Their meridians are a basis of first homology:
general-position bounding discs give generation, and the logarithmic
differentials of the distinct defining linear forms give the dual
independent classes. Its dimension is $b(b-3)/2$. The displayed $b$
relations follow by moving one point around all the other points on
the sphere; they are linearly independent over $\Q$, since a relation
among their rows has coefficients satisfying $a_i+a_j=0$ for all pairs.
The dimension count proves (E.13.4a).

The classes of distinct $e_{ij}$ are distinct when $b\ge5$. Otherwise
their difference would have coefficients $a_i+a_j$ supported on at
most two edges of the complete graph. The graph obtained by deleting
those edges is connected and contains an odd cycle; the zero
coefficients there force every $a_i=0$, a contradiction. A permutation
acting trivially on the quotient therefore fixes every two-element
subset, hence is the identity. The same proof applies after tensoring
with $\Q_\ell$, and so detects the permutation action on the
abelianization of the pure profinite group. A central element of any
coloured genus-zero group centralizes its pure subgroup. Its permutation
is consequently trivial, so it is in the centre-free pure subgroup
and is itself trivial. This proves the required centrelessness without
any assumption about the centres of the individual colour-permutation
groups.

The outer-triviality hypothesis permits gauges on $T,T_q$ making their
group automorphisms the identity while preserving the filling
$T_q\to T$. Here is the precise adjustment. First gauge the target
group action to the identity. The image of the filling is in its one
right orbit, so a source gauge fixes the filling. The remaining source
group action is inner, say by $h$. Compatibility with the filling makes
the image of $h$ central in $\widehat M_T$, hence trivial. Therefore
$h\in P_T$; a further source gauge by $h^{-1}$ fixes the filling and
makes the source action the identity. In particular $\psi_T$ in
(E.13.4) is inner, and may be taken to be the identity for that formula.

It follows at once from (E.13.4) that $f$ and $Ff$ have conjugate
image subgroups in $P_T$, hence the same unlabelled right cover orbit.
We next retain the necessary source labels as well.

\emph{Recovering a labelled simple ramification point.}
Let $c$ be one of the source punctures whose label is retained. Filling
just $c$ leaves an affine hyperbolic surface in the applications below.
By the deletion argument $\psi_S(N_c)=N_c$. Taking normal closures
in $P_T$ in (E.13.4) therefore gives
\[
 \overline{\langle\!\langle f_*(I_c)\rangle\!\rangle}_{P_T}
 =
 \overline{\langle\!\langle (Ff)_*(I_c)\rangle\!\rangle}_{P_T}.
                                                               \tag{E.13.5}
\]
For a covering arrow, including a completed one, a labelled source
cusp maps to a target cusp with a positive integral ramification index:
an actual covering has that property, and completed mapping classes
preserve the specified peripheral conjugacy classes by continuity.
If these target cusps and indices are $b_i,e$ for $f$ and $b_j,e'$
for $Ff$, the images of the two normal closures in
\[
P_T^{\rm ab}=\widehat\Z^{\,b}/
                 \widehat\Z(1,\ldots,1)
\]
are respectively $e\widehat\Z[c_{b_i}]$ and
$e'\widehat\Z[c_{b_j}]$. Since $b\ge3$, the indicated cusp lines
are pairwise distinct, including after tensoring with any $\Q_\ell$.
Equation (E.13.5) forces $i=j$. Equality of the two integral ideals
then gives $v_\ell(e)=v_\ell(e')$ for every prime $\ell$, and hence
$e=e'$. Thus target incidence and exact ramification index of every
retained source label are preserved. This uses normal closures and
genus-zero abelianization, not a profinite Magnus property.

In particular, if each retained source mark is the unique simple
ramification point above a separately distinguished target branch mark,
the isomorphism of the underlying covers necessarily respects those
source labels. Therefore $Ff$ and $f$ are in the same right orbit for
the required coloured source mapping-class group.

\emph{A Hurwitz space with connected source fibres.}
Fix a hyperbolic marked type $(g,n)$ with $g\ge1$, and choose
\[
                 d\ge2g+2n+5,\qquad b=2d+2g-2.
\]
Let $\mathcal H$ parametrize connected degree-$d$ simply branched
maps $C\to\mathbf P^1$, with $n$ distinct labelled simple
ramification points on $C$. Their branch values are separately labelled;
the other $b-n$ branch values are unordered. Targets are considered
modulo $\operatorname{PGL}_2$. A connected simply branched cover has
monodromy $S_d$: it is transitive and generated by transpositions.
For $d\ge3$ its deck group, the centralizer of this $S_d$ in its
degree-$d$ permutation representation, is trivial.

The usual Hurwitz construction gives a finite etale branch map
\[
 \mathcal H\longrightarrow\mathcal M_{0,n,[b-n]}.
                                                               \tag{E.13.6}
\]
The classical branch-cover construction is stated in
\citet[\S1.1, pp.~313--314]{E_RomagnyWewers2006}; the stack form
and non-Galois cover interpretation are the ones used in E.8
\citep[Proposition 6.14; Theorem 6.20]{E_BR2007}.
The stronger source-fibre assertion needed here is proved directly next,
rather than inferred from mere surjectivity onto $\mathcal M_g$.

Fix a pointed curve $(C;P_1,\ldots,P_n)$. Vary a reduced effective
divisor $D$ of degree $d$ avoiding the $P_i$. This parameter space is
a nonempty open subset of $\operatorname{Sym}^d(C)$, hence smooth
and irreducible. The spaces $H^0(C,\mathcal O_C(D))$ form a vector
bundle, since $d>2g-2$. Requiring
\[
                         df(P_i)=0\quad(1\le i\le n)
\]
gives $n$ independent linear conditions, by Riemann--Roch jet
surjectivity. Their kernels form a vector bundle on the divisor
parameter space. Inside it impose that every point of $D$ is an exact
simple pole, that the $P_i$ are simple ramification points, that every
other ramification point is simple, and that all branch values are
distinct.

These are nonempty open conditions. Here are the uniform dimension
checks, with jet lengths specified. Prescribing the value and first
derivative at every $P_i$ has total length $2n$. For the additional
conditions below, the largest needed finite subscheme has total length
$2n+4$: two further points contribute length two each; a multiple
critical point contributes length three; adding a second derivative
at a prescribed point increases the length $2n$ by one. The bound on
$d$ gives $H^1(C,\mathcal O_C(D-Z))=0$ for each such $Z$, and hence
surjectivity onto all these jets. A pole-coefficient condition is
handled likewise by $D-P$ at that pole, together with the length-$2n$
marked jets. Nonzero pole coefficient excludes a
proper hyperplane for each point of $D$. Failure of simple ramification
at a prescribed $P_i$ excludes another proper hyperplane. A multiple
critical point at an unspecified point imposes two independent linear
conditions over a one-dimensional locus. Two distinct unspecified
critical points with equal values impose three independent linear
conditions over a two-dimensional locus. If one critical point is a
prescribed $P_i$, the corresponding equal-value condition and the
other derivative condition impose two conditions over a one-dimensional
locus; equality at two prescribed critical points is one proper
hyperplane. Incidence dimension therefore shows that these bad loci
cannot cover the vector space. There is a finite union of these
incidence types. Simple poles are unramified over infinity, so no
additional critical point occurs at a pole.

Every function in this open set has unique pole divisor $D$.
Consequently the construction is exactly the coordinate chart of the
map space where infinity is unramified, not just a parametrizing
subset of unknown image. Its fibre over $(C;P_i)$ is irreducible.
The construction works relatively over $\mathcal M_{g,n}$:
the reduced relative symmetric power is smooth, the indicated
cohomology bundles are locally free, and the jet map has constant
surjective rank. Thus this chart is smooth and surjective over
$\mathcal M_{g,n}$, with geometrically connected fibres.

Every abstract target admits a coordinate with infinity outside its
finite branch locus. The choices form a smooth connected open subset
of $\operatorname{PGL}_2$. Passing from the chart to target coordinates
modulo $\operatorname{PGL}_2$ therefore shows that
\[
              \mathcal H\longrightarrow\mathcal M_{g,n}
                                                               \tag{E.13.7}
\]
is smooth and surjective with geometrically connected fibres. In
particular it induces a surjection on full etale fundamental groups.
Indeed every connected finite etale cover of the base pulls back to
a connected space: smooth surjectivity is open and the fibres are
connected, so a nontrivial clopen decomposition upstairs would give
one downstairs. This criterion is precisely surjectivity on the
finite-cover fundamental groups. The same argument shows that
$\mathcal H$ is connected.

The same construction works for every coloured or unordered version of
the $n$ retained source marks. The derivative conditions are the
vanishing of the differential on their finite etale marked divisor;
they do not require an ordering. The corresponding target branch marks
receive the same colours, while the other branch marks remain unordered.
One can verify the vector-bundle construction after an etale labelling
cover and then descend it; the geometric fibre is still the same space
of functions on the fixed curve, with no extra ordering choices.
Thus its connectedness and the fundamental-group surjection hold for
these versions as well.

\emph{Fixing the higher-genus action.}
Choose a point of $\mathcal H$. Remove all branch values from the
target, obtaining $T$, and remove all their inverse images from $C$,
obtaining $S$. The $n$ prescribed ramification points are distinguished
source punctures, and all other punctures are in the forgotten colour.
We have a cover $f:S\to T$ and a filling $u:S\to C\setminus\{P_i\}$.
The target has $b\ge5$ punctures and $b-n\ge3$, so the preceding
genus-zero normalization applies. The source has
$b(d-1)$ punctures, so every single deletion used to recover its
marked ramification profile remains affine hyperbolic.

By (E.13.4)--(E.13.5), $Ff$ lies in the same coloured right orbit
as $f$. For coloured retained marks, apply the labelled calculation
on the specified finite labelling cover and then forget the labels;
the resulting cover isomorphism respects each retained colour.
Gauge the source object to arrange $Ff=f$, keeping the
target action the identity. Let $H_f\subset\widehat M_T$ be the
target lifting subgroup. It is the fundamental group of the connected
Hurwitz cover (E.13.6), with the chosen basepoint. Since the deck
group is trivial, for $h\in H_f$ there is a unique lift
$\lambda_f(h)\in\widehat M_S$, characterized by
\[
                       h f=f\lambda_f(h).
\]
Applying $F$ and using the trivial right stabilizer of $f$ gives
\[
                     F_S(\lambda_f(h))=\lambda_f(h)
                     \qquad(h\in H_f).                       \tag{E.13.8}
\]

The image of the filling $u$ is in its unique right orbit. Write
$F(u)=u a$. The forgetting map
$\widehat M_S\to\widehat M_{g,n}$ is surjective, and
$u a=\operatorname{forget}(a)u$. A gauge on the filled object
therefore makes $F(u)=u$, without changing $Ff=f$ or (E.13.8).
Compatibility with $u$ then makes $F_{g,n}$ the identity on
\[
 \operatorname{im}\bigl(H_f\xrightarrow{\lambda_f}
                  \widehat M_S\longrightarrow\widehat M_{g,n}\bigr).
\]
Here left cancellation is valid for this filling: its right orbit is
the quotient by the kernel of the surjective forgetting homomorphism,
so the target group acts on that orbit by its regular left action.
This image is all of $\widehat M_{g,n}$ by (E.13.7). Hence
$F_{g,n}$ is the identity after gauge, so its original outer class is
trivial. The genus-zero cases were part of the hypothesis, and the
argument handles every remaining hyperbolic type, including proper
genus at least two and the once-punctured genus-one type.
\end{proof}

The conclusion is an injectivity statement for the complete genus-zero
outer-action system specified above. It does not identify every such
system with an element of $\widehat{\mathrm{GT}}$. An arithmetic-image
criterion phrased in that group must either prove the relevant comparison
or explicitly require that the complete genus-zero system is the standard
tower action of an element of the specified subgroup of
$\widehat{\mathrm{GT}}$. The theorem then supplies uniqueness of its
all-types outer extension whenever one already has an arithmetic lift.
Selecting the actual arithmetic subgroup from the conjugacy class in the
Hoshi--Mochizuki--Tsujimura theorem is a separate step. Neither a
genus-zero-to-$\widehat{\mathrm{GT}}$ comparison nor an assertion that
all genus-zero conjugators lift is proved here.
